\chapter{Approach}

Chương này trình bày chi tiết kiến trúc kỹ thuật của hệ thống "Multi-Agent Security Auditor". Khác với các hệ thống rà soát truyền thống hoạt động theo luồng tuyến tính (Linear Scripting), giải pháp của chúng tôi được xây dựng dựa trên kiến trúc **Graph-based Orchestration** (Điều phối dựa trên Đồ thị) sử dụng framework LangGraph. Kiến trúc này cho phép quản lý trạng thái phức tạp, hỗ trợ các quy trình có tính chu trình (cyclic) và đảm bảo khả năng mở rộng module hóa.

\section{Kiến trúc Điều phối Trung tâm (Graph-based Orchestration)}

Trung tâm của hệ thống là bộ điều phối `Graph Orchestrator`, hoạt động như một máy trạng thái hữu hạn (Finite State Machine). Cấu trúc này định nghĩa luồng dữ liệu di chuyển giữa các Tác nhân (Agents) thông qua một bộ nhớ chia sẻ.

\subsection{Quản lý Trạng thái (Shared State Management)}
Thay vì truyền dữ liệu trực tiếp giữa các hàm, hệ thống duy trì một đối tượng trạng thái toàn cục (Global State Object) xuyên suốt vòng đời của phiên kiểm toán. Cấu trúc dữ liệu này được định nghĩa như sau:

\begin{equation}
    \mathcal{S} = \{ R_{user}, \mathcal{P}_{scan}, \mathcal{J}_{ids}, \mathcal{F}_{raw}, \mathcal{F}_{enriched}, \mathcal{F}_{fail}, \text{Flag}_{fix} }
\end{equation}

Trong đó, dựa trên mã nguồn thực tế (\texttt{AgentState TypedDict}):
\begin{itemize}
    \item $\mathcal{P}_{scan}$: Kế hoạch quét được sinh ra bởi Planning Agent (JSON Config).
    \item $\mathcal{J}_{ids}$: Danh sách Job ID của các tiến trình Prowler đang chạy ngầm.
    \item $\mathcal{F}_{enriched}$: Danh sách các phát hiện (Findings) đã được Risk Agent làm giàu thông tin và chấm điểm.
    \item $\mathcal{F}_{fail}$: Tập hợp con các phát hiện có trạng thái "FAIL" cần khắc phục.
    \item $\text{Flag}_{fix}$: Cờ trạng thái (\texttt{remediation\_needed}) xác định xem hệ thống có cần kích hoạt quy trình sửa lỗi hay không.
\end{itemize}

\subsection{Sơ đồ Luồng xử lý (Workflow Pipeline)}
Hệ thống được mô hình hóa thành một đồ thị có hướng (Directed Graph) với các đỉnh (Nodes) là các Agent chuyên biệt. Luồng thực thi tuân theo quy trình khép kín:

\begin{enumerate}
    \item \textbf{Plan Node:} Phân tích yêu cầu người dùng $\rightarrow$ Tạo cấu hình quét.
    \item \textbf{Scan Node:} Điều phối Scanner Agent gọi tool Prowler.
    \item \textbf{Monitor Node:} Giám sát tiến trình và thu thập dữ liệu thô (Raw Findings).
    \item \textbf{Risk Eval Node:} Đánh giá mức độ nghiêm trọng của từng lỗi.
    \item \textbf{Assessment Node:} Lọc ra các lỗi vi phạm ngưỡng an toàn (Failures).
    \item \textbf{Human Review Node:} (Điểm ngắt) Chờ xác nhận từ người dùng trước khi can thiệp hạ tầng.
    \item \textbf{Remediate Node:} (Có điều kiện) Thực thi sửa lỗi nếu được duyệt.
    \item \textbf{Rescan \& Analysis Node:} Quét lại để kiểm chứng và so sánh sự thay đổi.
    \item \textbf{Report Node:} Tổng hợp dữ liệu trạng thái cuối cùng để xuất báo cáo.
\end{enumerate}


\section{Cơ chế Hoạt động của các Tác nhân Chuyên biệt}

Hệ thống áp dụng nguyên lý "Chuyên môn hóa Tác nhân", trong đó mỗi Agent sử dụng một chiến lược Prompting và cơ chế gọi Tool (Tool Usage) riêng biệt.

\subsection{Scanner Agent: Dynamic Tool Binding}
\textit{Scanner Agent} chịu trách nhiệm tương tác với công cụ rà soát Prowler. Điểm cải tiến trong phiên bản này là cơ chế **Dynamic Tool Binding** thay vì Hardcoding.

\begin{itemize}
    \item \textbf{Khởi tạo:} Agent tự động nạp danh sách công cụ từ module \texttt{agent\_tools} và tạo ra một bảng ánh xạ động (\texttt{tools\_map}).
    \item \textbf{Cơ chế gọi:} Sử dụng phương thức \texttt{bind\_tools} của LangChain để gắn kèm định nghĩa hàm vào LLM. Khi nhận nhiệm vụ, Llama 3.2 sẽ tự động suy luận và trả về cấu trúc \texttt{tool\_calls} chứa tên tool và tham số cần gọi.
    \item \textbf{Thực thi:} Agent tra cứu tên tool trong bảng ánh xạ và thực thi hàm tương ứng (ví dụ: \texttt{run\_prowler\_scan}), sau đó trích xuất Job ID để theo dõi.
\end{itemize}

\subsection{Risk Evaluation Agent: JSON Mode \& Scoring Rubric}
Đây là tác nhân phân tích, chịu trách nhiệm chuyển đổi dữ liệu thô thành thông tin tình báo (Intelligence).

\begin{itemize}
    \item \textbf{Input Processing:} Agent lọc danh sách các finding có trạng thái \texttt{FAIL} và loại bỏ các trường thông tin nhiễu.
    \item \textbf{Structured Output (JSON Mode):} Agent ép buộc mô hình ngôn ngữ trả về định dạng JSON nghiêm ngặt bằng tham số \texttt{format="json"}. Điều này đảm bảo tính nhất quán cho các bước xử lý máy học phía sau.
    \item \textbf{Scoring Rubric (Thang chấm điểm):} System Prompt chứa một bộ quy tắc chấm điểm rủi ro chi tiết (Rubric):
    \begin{itemize}
        \item \textit{Critical (9-10):} Lỗi Public Access, lộ Key.
        \item \textit{High (7-8):} Cấu hình sai nghiêm trọng nhưng giới hạn trong nội bộ.
        \item \textit{Medium (4-6):} Thiếu Logging, Monitor.
        \item \textit{Low (1-3):} Các lỗi về Tagging hoặc thông tin.
    \end{itemize}
    \item \textbf{Nhiệm vụ:} Với mỗi finding, Agent sẽ đọc \texttt{description} và \texttt{status\_details} để gán nhãn \texttt{severity} và \texttt{risk\_score} chính xác.
\end{itemize}

\subsection{Remediate Agent: Semantic Tool Selection}
Tác nhân khắc phục sự cố được thiết kế theo hướng **Fully Dynamic Inference**, cho phép hệ thống tự động tìm công cụ sửa lỗi phù hợp mà không cần lập trình trước các cặp "Lỗi - Cách sửa".

Quy trình suy luận diễn ra như sau:
\begin{enumerate}
    \item \textbf{Context Loading:} Agent nạp toàn bộ danh sách mô tả của các công cụ sửa lỗi (ví dụ: \texttt{s3\_enable\_versioning}, \texttt{iam\_rotate\_key}) vào System Prompt.
    \item \textbf{Semantic Matching:} Khi nhận được một lỗi (ví dụ: "S3 Bucket X thiếu Versioning"), LLM sẽ thực hiện so sánh ngữ nghĩa:
    \begin{equation}
        \text{Tool}_{selected} = \text{argmax}_{t \in \mathcal{T}} (\text{Similarity}(\text{Error\_Desc}, t_{desc}))
    \end{equation}
    Mô hình sẽ tự động hiểu rằng từ khóa "Versioning" trong lỗi khớp với công cụ \texttt{s3\_enable\_versioning}.
    \item \textbf{Safety Guardrails:} Agent bao gồm logic kiểm tra an toàn:
    \begin{itemize}
        \item Tự động bỏ qua nếu \texttt{resource\_id} là Account ID (để tránh tác động diện rộng).
        \item Chỉ thực thi khi độ tự tin (Confidence) của mô hình cao thông qua việc so khớp chính xác tên tool.
    \end{itemize}
\end{enumerate}

\section{Quản lý Dữ liệu và Tính Bền vững (Data Persistence)}

Do tính chất chạy cục bộ (Local Execution), hệ thống sử dụng chiến lược lưu trữ dựa trên tệp (File-based Storage) để đảm bảo hiệu năng cao và dễ dàng kiểm tra (audit).

\begin{itemize}
    \item \textbf{Snapshots:} Tại các điểm chuyển giao quan trọng của đồ thị (Graph Nodes), dữ liệu được lưu xuống ổ cứng dưới dạng JSON:
    \begin{itemize}
        \item \texttt{data/initial\_scan\_config.json}: Lưu cấu hình phạm vi quét.
        \item \texttt{data/pre\_scan.json}: Lưu trạng thái hệ thống trước khi can thiệp.
        \item \texttt{data/analysis\_diff.json}: Lưu kết quả so sánh (Diff) giữa trước và sau khi sửa lỗi.
    \end{itemize}
    \item \textbf{Report Generation:} \textit{Report Agent} đọc dữ liệu từ các tệp snapshot này để tổng hợp thành báo cáo Markdown cuối cùng, đảm bảo rằng báo cáo luôn phản ánh đúng dữ liệu thực tế đã xử lý.
\end{itemize}